\section{Verification}\label{sec:verify}

{\sloppy
This section says how the numbers in this paper were computed and which
results are also proved in Lean. All computations are in Python. The values of $R_N$ in Table~\ref{tab:heavy}
were computed in two ways, both from Proposition~\ref{prop:tworenewal},
namely by its formula and by a spectral method that reads each factor off as
a Taylor coefficient of a generating function. For small $N$ they were also
checked against a recursion on the first run and a sum over all run
compositions, which do not use Proposition~\ref{prop:tworenewal}. The aging
probabilities, and with them the crossing $c_{200}$, were computed in three
ways, namely by a forward recursion, from transfer matrices at roots of unity,
and from the sums over paths in~\eqref{eq:aging-ratio}. In exact rational
arithmetic $C_{200}-E_{200}$ is negative at $c_{200}-10^{-6}$ and positive at
$c_{200}+10^{-6}$.

No proof in this paper rests on a computation. The tables,
the figures and the numbers quoted outside the proofs rest on computation
alone. The script \texttt{verify\_heavy\_runs.py} recomputes the numbers of
Section~\ref{sec:heavy} and Appendix~\ref{app:heavy}, the script \texttt{verify\_aging.py}
those of Section~\ref{sec:aging} and Appendix~\ref{app:aging}, and
\texttt{verify\_aging\_urn.py} checks Proposition~\ref{prop:aging-urn} in exact arithmetic for
two, three and four states. The numbers in the caption of
Figure~\ref{fig:aging} come from \texttt{figures/make\_aging\_figure.py}.

Many statements of this paper are also proved in Lean~4.33.1 with Mathlib at commit
\texttt{0df444a360ea}, in the folder \nolinkurl{kagey-problems/problem131/lean/PaperD}, which
is shared with an earlier version of this work. The file \nolinkurl{MAP.md} there lists each
statement of this paper by its label, with its Lean names, its status in Lean and the source of
each classical input. The kernel checks every proof, and every theorem depends only on the
axioms \texttt{propext}, \texttt{Classical.choice} and \texttt{Quot.sound}.
The statement-by-label table in \nolinkurl{MAP.md} and the detailed table
\nolinkurl{coverage_F.tex} record the coverage and its qualifications. Some proofs use six classical
results, which Lean states as named hypotheses, and nothing else is assumed. These are the local
limit theorem of Gnedenko (GN)~\cite[Theorem~4.2.1]{ibragimovlinnik1971}, the expansion of
$\mathrm{Li}_s(e^{-v})$ near $v=0$ (POLY)~\cite[25.12.12]{dlmf}, the continuity theorem of
L\'evy (LEVY) and Fourier inversion for an integrable characteristic function
(FI)~\cite[Section~XV.3]{feller1971}, the continuity theorem for Laplace transforms
(LAP)~\cite[Section~XIII.1]{feller1971} and the uniqueness theorem of Lerch for Laplace
transforms (LERCH)~\cite[Chapter~II]{widder1941}.



One proof in Lean differs from ours. In Step~4 of the proof of Theorem~\ref{thm:lamperti},
Lean finds $\int_0^\infty tg(t)^2\,dt$ from $\E G^{i\tau}$ by Parseval's identity for the
density of $\log G$, so it does not identify the law of $G_1/(G_1+G_2)$. Lean does not
formalize the six classical results, Theorem~\ref{thm:alpha2},
Remarks~\ref{rem:humps}, \ref{rem:aging-feller} and~\ref{rem:other}, or
the numbers that rest on computation.
\par}

\paragraph{Open questions.}
For the heavy-tailed walk, a rate in Theorem~\ref{thm:center} remains open.
So does monotonicity of $E_N/C_N$ in $N$, which would give uniqueness of a
crossover size when the two probabilities exchange order.
